\scanpage{17}
% L01–03
Cas c)

\deleted{d'où $H(k)\simeq\gamma(k)\,\unsure{\ltimes|\cdot}\,U(k)$ $(U(k)\simeq k\text{ ou }k^*)$ et}

% L04–09
$(\ast)\Longrightarrow(\ast a)$ \deleted{si $p$ impair}, car les ordres des élts d'ordre fini de $\mathrm{GP}(1)$ sont les $i$ premiers à $p$, \deleted{et $p$ \illegible{1}} (correspondant aux él. $\frac12$ simples) et $p$ (correspondant aux éléments unipotents). \deleted{Mais \inserted{et} si $i=p=2$, \unsure{$G$} $\mathbb D_i=\mathbb Z_2\times\mathbb Z_2$ n'est pas contenu dans $\mathrm{GP}(1)$ (car \unsure{les uni} les $G_a$ \inserted{$\subset\mathrm{GP}(1)$} sont les unipotents maximaux). Et}

% L10–12
$(\ast a)\Longrightarrow(\ast)$ immédiat, en distinguant et examinant (pour le cas $i=p$) le cas $p$ impair, et le cas $p$ pair (alors cas chevauchant, l'image de $\mathbb D_2$ est à choisir dans $U$ !).

% L13–19
\deleted{\uline{Cor.} \unsure{Soit $p$, $i$ préc.|Soit $i$ pair} \unsure{Supposons $i$ pair.} Considérons}
\[
\deleted{(\ast')\quad\exists\ \text{repr. fidèle}\quad\mathbb D_i\times\mathbb Z_{/2}\longrightarrow\mathrm{GP}(1,k).}
\]
\deleted{a) Si $p=0$, \deleted{ou $i=1$} $(\ast')$ est satisfaite sss $i=0$}

\deleted{b) Si $i=0$, \inserted{donc $\mathbb D_i\times\mathbb Z_{/2}=\deleted{\unsure{\mathbb D_i}}=\mathbb D_\infty\times\mathbb Z_{/2}$}, $(\ast')$ est satisfaite sss \deleted{\unsure{$p$ est pair}} $p\ne2$.}

\deleted{c) Si $i\geq2$, $p>0$, alors $(\ast')$ n'est pas satisfaite.}

\deleted{\uline{Dém} \unsure{a) immédiat} (crit. de parité}

% L20–26
On s'intéresse aux repr. de $\mathbb D_\infty$ \inserted{$=\unsure{\mathrm G_{\Delta_1}|G_{P1}|G_1}=\{\underline\tau_0,\underline\tau_1\mid\underline\tau_0^2=\underline\tau_1^2=1\}$}, dans des $G$ \inserted{$\subset\uline{\mathrm{Aut}}(X)$} formes de $\mathrm{GP}(1)$ (sur $S$) \deleted{$\Longleftrightarrow$} la cat de celles-ci est équivalente à celle des représentations \deleted{\unsure{en}} projectives de dim $2$ de $\mathrm G_1$, \inserted{(dans un fibré projectif $X$)} \uline{munies de} \deleted{repr.} quelques bons invariants.
% slanted
\margin{On pose $\tau_i=\varphi(\underline\tau_i)$ $(i=0,1)$ : $u=\tau_0\tau_1=\varphi(\underline u)$}

% L27–31
\deleted{On les} \deleted{On va regarder} on va regarder d'abord le cas où la condition suivante est satisfaite

\deleted{$(\ast)$ $\forall s\in S$, $u_s$ $(=\tau_{0s}\tau_{1s})$ n'est}

\[
\begin{aligned}
(\ast)\quad &\forall s\in S,\qquad u_s^2\ne1\qquad
(\text{i.e. }\varphi_s:\mathbb D_\infty\longrightarrow G_s\ \text{ne se factorise}\\
&\hspace{18em}\text{pas par }\mathbb D_2)\\
&\text{i.e. }\tau_{0s},\tau_{1s}\text{ ne \uline{commutent pas}}
\end{aligned}
\]

% L32–35
\deleted{La catégorie de ces \unsure{fon} couples $(G,\rho)$ s'identifie à}

Or \deleted{\unsure{on sait}}, si on introduit $a_0,H_0,a_1,H_1$ comme \deleted{à la cat. des polygônes réguliers} d'habitude, la condition $(\ast)$ implique
\[
a_{0s}\notin H_1\quad(\Longleftrightarrow{}a_1\notin H_0)
\]
\inserted{\unsure{pour : le G inv.}} en tte fibre
% archive sheet number: 16
